\begin{table}[t]
\centering
\begin{threeparttable}
\caption{Direct test of the strategy-vs-outcome dissociation in privileged increment (analysis population).}
\label{tab:tab-dissociation-216}
\small
\begin{tabular}{lrrrrrr}
\toprule
limiting contrast & leak contrast & $\Delta$ limiting & $\Delta$ leak & difference & 95\% CI & p \\
\midrule
limiting vs direct (disclosers only) & substantive leak vs appropriate & +0.162 & +0.057 & +0.104 & [-0.017, 0.220] & .045 \\
limiting vs direct (disclosers only) & substantive leak vs rest & +0.162 & +0.028 & +0.130 & [0.023, 0.242] & .010 \\
limiting vs direct (all) & substantive leak vs appropriate & +0.125 & +0.057 & +0.069 & [-0.055, 0.181] & .144 \\
limiting vs direct (all) & substantive leak vs rest & +0.125 & +0.028 & +0.096 & [-0.010, 0.207] & .042 \\
\bottomrule
\end{tabular}
\begin{tablenotes}[flushleft]\footnotesize
\item A difference between a significant and a non-significant result is not itself significant (Gelman \& Stern 2006, The American Statistician 60(4)). This table tests the difference of the two deltas directly.
\item Draws are aligned by index across contrasts (same seed and draw order), so the comparison is paired by construction.
\item The comparison against substantive\_leak is supported; the comparison against leak\_vs\_appropriate has an interval that grazes zero and is reported as suggestive. No multiplicity adjustment is applied across these four comparisons --- they are reported as a coherent set, not as independent tests.
\end{tablenotes}
\end{threeparttable}
\end{table}
